% BIB OAI-Preprint: {AUTHOR}, \emph{{PREPRINT TITLE}}, {RELEASE NAME} preprint, {DATE}, \url{{RELEASE URL, e.g. https://github.com/openai/math}}.
% BIB OAI-Companion1: {AUTHOR}, \emph{{COMPANION TITLE}}, {RELEASE NAME} preprint, {DATE} (cited by the preprint as \texttt{{BIB KEY}}).
\section{About this document}\label{sec:about}

\paragraph{The preprint.}
The preprint is \emph{{PREPRINT TITLE}}, dated {DATE}, with author line
``{AUTHOR}'' \cite{OAI-Preprint}. {ONE SENTENCE OF CONTEXT, e.g. the release it
belongs to.} It proves {ONE OR TWO SENTENCES: THE MAIN RESULT}.
Section~\ref{sec:overview} summarizes the paper and fixes the notation used
throughout this document. {EITHER: The preprint relies on {N} companion
preprints from the same release, \cite{OAI-Companion1}, which it cites as
\texttt{{BIB KEY}}. They belong to the same release, so they are not treated as
earlier work here. OR: The preprint cites no other preprint from the same
release.}

\paragraph{What this document reports.}
Each finding is a citation the preprint is missing: part of the preprint's
mathematics (a technique, construction, estimate, proof route or result) that
an identifiable earlier paper already contains, where the preprint does not
cite that paper for it where it is used. ``Earlier'' means publicly available
before the preprint's date. The one exception, overstated novelty, is
described below. Not reported: textbook material in its textbook form; weaker versions of
intermediate statements; and, for content the preprint uses, cases where it
cites a later source that covers that content rather than the original. An
earlier result on the same problem is reported even when later work that the
preprint cites implies it. Findings are ordered by centrality, then by relationship in
the order listed below, then by position in the preprint.

\paragraph{Relationships.}
A finding's relationship is the one between the preprint and the first
earlier work listed for it, the work with the highest relationship in the
order below. Any other works listed relate to the same content more weakly or
in part; the finding says how, including which work originated the idea when
that is not the first.
\begin{description}[leftmargin=1.5em,font=\normalfont\itshape]
\item[same result.] An earlier paper, or several earlier papers together, prove essentially the same statement.
\item[same technique.] The earlier technique, construction or theory is
reused essentially intact, possibly in a new setting or with new ingredients
around it.
\item[adapted.] The earlier technique's structure is followed, with a change
at its core: a different object, or a key step replaced.
\item[proof shortcut.] A published result would replace or substantially
shorten a large part of the argument.
\item[same route.] An earlier paper proposed the proof strategy, even
heuristically, and the preprint carries it out.
\item[fragment.] One identifiable piece of the earlier technique is reused.
\item[weaker result.] An earlier paper proves a weaker, special-case or
partial-range version of one of the preprint's main results, or of a
statement it presents as a significant contribution.
\end{description}

\paragraph{Centrality.}
Each finding also gets a best-effort judgment of how much of the preprint's
main argument depends on it: \emph{main result} (reserved for earlier work that
covers the main result itself, or most of it), \emph{central},
\emph{supporting} or \emph{peripheral}. An earlier weaker result is never
labelled main result. Each finding states the chain of results it feeds, so
the basis for the judgment is visible.

\paragraph{Citations.}
The ``Cited?'' column of the summary names the earlier work and says what the
preprint cites for the content in question: \emph{not cited} (the work is not
in the bibliography), \emph{cited for other content}, \emph{credited elsewhere}
(cited for this content, but not where it is used,
in a way that an expert reader who doesn't know the earlier work might not connect with it), or \emph{credited in
part}. A citation counts only when it is \emph{for that specific content};
citing the same paper nearby for a different result does not count. Content
the preprint credits where it is used, directly or through a cited location
that itself credits the earlier work{IF COMPANIONS: (such as a companion
lemma)}, is not reported, except when its novelty is overstated (below). For
overstated novelty, the column says where the preprint does cite the work.

\paragraph{Flags.}
\begin{description}[leftmargin=1.5em,font=\normalfont\bfseries]
\item[Attribution gap.] Marks the most serious missing citations: the
relationship is same technique, adapted or same result; the centrality is main
result, central or supporting; and the preprint cites the principal earlier
work for this content nowhere.
\item[Overstated novelty.] The preprint does cite the earlier work for this
content, but its abstract, a statement of its significant contributions, or the framing of
a main result presents that content as new, explicitly or in the way a reader
would reasonably take it. This is the one case in which content the preprint
credits is reported.
\end{description}

\paragraph{How to check a claim.}
Locations in the preprint{IF COMPANIONS: and its companions} use its section,
equation and theorem numbers and the page numbers printed in its PDF.
Locations in other sources give the theorem, equation and printed page of the
version listed in the bibliography. Each finding ends with a short checklist
of the passages to put side by side.

% Provenance: keep this block commented out until the publish step, then
% remove the leading "% " from each line and fill in the placeholders. It
% names the user, their models and their input: publish only the exact text
% the user approved (references/publish.md, section 1).
% \paragraph{Provenance.}
% This document was generated by LLM agents ({MODEL}) running the
% precedence-review skill
% (\href{https://github.com/jnshi/math-lineage/tree/{COMMIT}/.claude/skills/precedence-review}{jnshi/math-lineage},
% commit \texttt{{COMMIT}}), makes no claim of completeness, and has not been
% vetted by a human. It was run by {NAME}{IF THEY SUPPLIED INPUT THAT SHAPED A
% FINDING: , who also supplied the leads behind finding~\ref{f:<finding>}}.{ANY
% DISCLOSURE THE OPERATOR ASKS TO ADD, VERBATIM.}
